\documentclass[11pt]{article}
\usepackage[a4paper,margin=18mm]{geometry}
\usepackage{fontspec}
\setmainfont{Latin Modern Roman}
\usepackage{amsmath,amssymb,amsthm,xfrac,xspace,hyperref}
\newcommand{\ie}{i.e.,\xspace}
\newcommand{\eg}{e.g.,\xspace}
\newcommand{\etc}{etc.\xspace}
\newcommand{\cf}{cf.\xspace}
\newcommand{\Subsection}{\subsection}
\newcommand{\Subsubsection}{\subsubsection}
\newcommand{\skpt}{\smallskip}
\newcommand{\smaller}{\small}
\newcommand{\backmatter}{}
\input{author-preamble.tex}
\emergencystretch=3em
\begin{document}
\begin{center}{\Large Stable item 1921}\end{center}
\paragraph{Source and attribution.}
David Constantine, Jean-François Lafont, Daniel J. Thompson, \emph{Strong symbolic dynamics for geodesic flows on CAT\$(-1)\$ spaces and other metric Anosov flows}, Journal de l’École polytechnique — Mathématiques 7 (2020), publisher version, DOI 10.5802/jep.115. \href{https://jep.centre-mersenne.org/articles/10.5802/jep.115/}{Publisher article}. Authors' text: \href{https://creativecommons.org/licenses/by/4.0/}{CC BY 4.0}.
\paragraph{Editorial source boundary.}
Prove Theorem A below for the original fixed-point-free Hölder metric Anosov flow with the additional explicit availability of the stated Hölder pre-Markov proper families at arbitrarily small positive scales. This is the scale condition actually used in the printed proof in Section 3.4. The selected conclusion is strong Markov coding: a finite-alphabet suspension with positive Hölder roof and a finite-to-one continuous surjective flow semiconjugacy, injective on a residual set, whose coding map is Hölder. No irreducibility is claimed without transitivity. Theorem A's proof uses the printed definitions and classical Bowen--Pollicott results quoted with references; these are retained as standard prerequisites, not independently reproved. The article's CAT(-1) geometric applications and further statistical conclusions are context only, not additional selected results.
\paragraph{Difficulty (editorial): H2.}
Metric hyperbolic orbit tracking is transferred through Hölder flow-box projections to exponential control of the symbolic coding map. Refinement of pre-Markov rectangles preserves return-time regularity, yielding a Hölder suspension roof and finite-to-one equivariant coding beyond smooth Anosov manifolds.
\paragraph{Transcription note (editorial).}
Author mathematical text and ordering are unchanged. Only publisher front matter is replaced by this independent wrapper. Original TeX body and macros remain separate editable inputs. Bibliography entries are recovered from the same cached publisher article HTML. No mathematical repair or assistant-generated proof is supplied.
\clearpage
\input{author-body.tex}
\begin{thebibliography}{99}
\bibitem{BPP16} [BAPP19] A. Broise-Alamichel, J. Parkkonen \& F. Paulin - Equidistribution and counting under equilibrium states in negative curvature and trees , Progress in Math. , vol. 329 , Birkhauser, 2019
\bibitem{bcls} [BCLS15] M. Bridgeman, R. Canary, F. Labourie \& A. Sambarino - “The pressure metric for Anosov representations” , Geom. Funct. Anal. 25 (2015), p. 1089-1179
\bibitem{BCLS18} [BCLS18] M. Bridgeman, R. Canary, F. Labourie \& A. Sambarino - “Simple root flows for Hitchin representations” , Geom. Dedicata 192 (2018) no. 1, p. 57-86
\bibitem{BC17} [BCS17] M. Bridgeman, R. Canary \& A. Sambarino - “An introduction to pressure metrics for higher Teichmüller spaces” , Ergodic Theory Dynam. Systems (2017), p. 1-35
\bibitem{bh} [BH99] M. R. Bridson \& A. Haefliger - Metric spaces of non-positive curvature , Grundlehren Math. Wiss. , vol. 319 , Springer, Berlin, 1999
\bibitem{mB95} [Bou95] M. Bourdon - “Structure conforme au bord et flot géodésique d’un CAT(-1)-espace” , Enseign. Math. 41 (1995) no. 1-2, p. 63-102
\bibitem{bowen-periodic} [Bow72] R. Bowen - “Periodic orbits for hyperbolic flows” , Amer. J. Math. 94 (1972) no. 1, p. 1-30
\bibitem{bowen-symbolic} [Bow73] R. Bowen - “Symbolic dynamics for hyperbolic flows” , Amer. J. Math. 95 (1973) no. 2, p. 429-460
\bibitem{bowen-ruelle} [BR75] R. Bowen \& D. Ruelle - “The ergodic theory of Axiom A flows” , Invent. Math. 29 (1975) no. 3, p. 181-202
\bibitem{BW} [BW72] R. Bowen \& P. Walters - “Expansive one-parameter flows” , J. Differential Equations 12 (1972), p. 180-193
\bibitem{Cannon} [Can84] J. W. Cannon - “The combinatorial structure of cocompact discrete hyperbolic groups” , Geom. Dedicata 16 (1984), p. 123-148
\bibitem{champetier} [Cha94] C. Champetier - “Petite simplification dans les groupes hyperboliques” , Ann. Fac. Sci. Toulouse Math. (6) 3 (1994) no. 2, p. 161-221
\bibitem{CLT} [CLT19] D. Constantine, J.-F. Lafont \& D. Thompson - “The weak specification property for geodesic flows on CAT(-1) spaces” , 2019, to appear in Groups Geom. Dyn.
\bibitem{cp} [CP12] M. Coornaert \& A. Papadopoulos - “Symbolic coding for the geodesic flow associated to a word hyperbolic group” , Manuscripta Math. 109 (2012), p. 465-492
\bibitem{DP84} [DP84] M. Denker \& W. Philipp - “Approximation by Brownian motion for Gibbs measures and flows under a function” , Ergodic Theory Dynam. Systems 4 (1984) no. 4, p. 541-552
\bibitem{das-simmons-urbanski} [DSU17] T. Das, D. Simmons \& M. Urbański - Geometry and dynamics in Gromov hyperbolic metric spaces , Math. Surveys and Monographs , vol. 218 , American Mathematical Society, Providence, RI, 2017
\bibitem{FH} [FH18] T. Fisher \& B. Hasselblatt - Hyperbolic Flows , Lect. Notes in Math. Sciences , vol. 16 , University of Tokyo, 2018, available at https://www.ms.u-tokyo.ac.jp/publication\_e/lecturenote\_e.html
\bibitem{GM} [GM10] K. Gelfert \& A. E. Motter - “(Non)invariance of dynamical quantities for orbit equivalent flows” , Comm. Math. Phys. 300 (2010) no. 2, p. 411-433
\bibitem{gromov} [Gro87] M. Gromov - “Hyperbolic groups” , in Essays in group theory (S. Gersten, ed.) , MSRI Publications , vol. 8 , Springer, 1987, p. 75-265
\bibitem{KT17} [KT19] T. Kucherenko \& D. J. Thompson - “Measures of maximal entropy for suspension flows over the full shift” , 2019, to appear in Math. Z.
\bibitem{mineyev} [Min05] I. Mineyev - “Flows and joins of metric spaces” , Geom. Topol. 9 (2005), p. 403-482
\bibitem{MT04} [MT04] I. Melbourne \& A. Török - “Statistical limit theorems for suspension flows” , Israel J. Math. 144 (2004), p. 191-209
\bibitem{pollicott} [Pol87] M. Pollicott - “Symbolic dynamics for Smale flows” , Amer. J. Math. 109 (1987) no. 1, p. 183-200
\bibitem{PP} [PP90] W. Parry \& M. Pollicott - Zeta functions and the periodic orbit structure of hyperbolic dynamics , Astérisque , vol. 187-188 , Société Mathématique de France, Paris, 1990
\bibitem{PS17} [PS17] R. Potrie \& A. Sambarino - “Eigenvalues and entropy of a Hitchin representation” , Invent. Math. 209 (2017) no. 3, p. 885-925
\bibitem{iP2014} [Put14] I. F. Putnam - A homology theory for Smale spaces , Mem. Amer. Math. Soc. , vol. 232, no. 1094 , American Mathematical Society, Providence, RI, 2014
\bibitem{ratner} [Rat73] M. Ratner - “The central limit theorem for geodesic flows on n -dimensional manifolds of negative curvature” , Israel J. Math. 16 (1973), p. 181-197
\bibitem{ratner2} [Rat74] M. Ratner - “Anosov flows with Gibbs measures are also Bernoullian” , Israel J. Math. 17 (1974), p. 380-391
\bibitem{ricks} [Ric17] R. Ricks - “Flat strips, Bowen-Margulis measures, and mixing of the geodesic flow for rank one CAT(0) spaces” , Ergodic Theory Dynam. Systems 37 (2017), p. 939-970
\bibitem{Roblin} [Rob03] T. Roblin - Ergodicité et équidistribution en courbure négative , vol. 95 , Société Mathématique de France, Paris, 2003
\bibitem{Ruelle} [Rue76] D. Ruelle - “A measure associated with Axiom-A attractors” , Amer. J. Math. 98 (1976) no. 3, p. 619-654
\bibitem{aS16} [Sam16] A. Sambarino - “On entropy, regularity and rigidity for convex representations of hyperbolic manifolds” , Math. Ann. 364 (2016) no. 1-2, p. 453-483
\bibitem{tapie} [Tap11] S. Tapie - “A variation formula for the topological entropy of convex-cocompact manifolds” , Ergodic Theory Dynam. Systems 31 (2011), p. 1849-1864
\bibitem{troyanov} [Tro90] M. Troyanov - “Espaces à courbure négative et groupes hyperboliques” (E. Ghys \& P. de la Harpe, eds.) , Progress in Math. , vol. 83 , Birkhäuser, 1990, p. 47-66
\end{thebibliography}
\end{document}
